% concatenated from Geometry_of_Lim_s_Condition.tex
\documentclass[11pt]{amsart}
\usepackage{xcolor}
\usepackage{cancel}
\usepackage{setspace}
\renewcommand{\baselinestretch}{2.25}
\usepackage{lscape}
\newcommand{\cL}{\mathcal L}
\newcommand{\F}{\mathcal{F}}
\newcommand{\I}{\mathcal {I}}
\newcommand{\cA}{\mathcal A}
\newcommand{\cD}{\mathcal D}
\newcommand{\cS}{\mathcal S}
\newcommand{\mb}{\makebox}
\newcommand{\es}{\emptyset}
\newcommand{\ci}{\subseteq}
\renewcommand{\u}{\cup}
\renewcommand{\i}{\cap}
\newcommand{\ra}{\rightarrow}
\newcommand{\Ra}{\Rightarrow}
\newcommand{\Lra}{\Leftrightarrow}
\newcommand{\lgra}{\longrightarrow}
\newcommand{\Lgra}{\Longrightarrow}
\newcommand{\al}{\alpha}
\newcommand{\be}{\beta}
\newcommand{\de}{\delta}
\newcommand{\e}{\varepsilon}
\newcommand{\s}{\sigma}
\newcommand{\la}{\lambda}
\newcommand{\pf}{\noindent {\bf Proof :} }
\newcommand{\ov}{\overline}
\newcommand{\ol}{\overline}
\newcommand{\iy}{\infty}
\newcommand{\N}{\mathbb{N}}
\newcommand{\K}{\mathcal{K}}
\newcounter{cnt1}
\newcounter{cnt2}
\newcommand{\blr}{\begin{list}{$(\roman{cnt1})$}
 {\usecounter{cnt1} \setlength{\topsep}{0pt}
 \setlength{\itemsep}{0pt}}}
\newcommand{\bla}{\begin{list}{$($\alph{cnt2}$)$}
 {\usecounter{cnt2} \setlength{\topsep}{0pt}
 \setlength{\itemsep}{0pt}}}
\newcommand{\el}{\end{list}}
\newtheorem{thm}{Theorem}
\newtheorem{lem}[thm]{Lemma}
\newtheorem{cor}[thm]{Corollary}
\newtheorem{ex}[thm]{Example}
\newtheorem{Q}[thm]{Question}
\newtheorem{Def}[thm]{Definition}
\newtheorem{note}[thm]{Note}
\newtheorem{prop}[thm]{Proposition}
\newtheorem{rem}[thm]{Remark}
\newcommand{\Rem}{\begin{rem} \rm}
\newcommand{\bdfn}{\begin{Def} \rm}
\newcommand{\edfn}{\end{Def}}
\renewcommand{\baselinestretch}{1.25}
\newcommand{\TFAE}{the following are equivalent~: }
\sloppy
\begin{document}
\large
\title[Geometry]{Some Geometric Aspects Related to Lim's Condition}
\author[Gothwal]{Deepak Gothwal}
\address[Deepak Gothwal]
{Department of Mathematics\\
Indian Institute of Technology, Kharagpur\\
West Bengal-721302\\ India,
\textit{E-mail~:}
\textit{deepakgothwal190496@gmail.com}}
\author[Rao]{T. S. S. R. K. Rao}
\address[T. S. S. R. K. Rao]
{Department of Mathematics\\
Shiv Nadar Institution of Eminence \\
Gautam Buddha Nagar\\ UP-201314 \\ India,
\textit{E-mail~:}
\textit{srin@fulbrightmail.org}}
\subjclass[2000]{Primary 46 B20, 46B22, 46B25, 46B50, 46E15, 46L05, 46L10, 47H10, 47L05, 47L15}
\keywords{Lim's condition, property ($\ddag$), $w^*$-normal structure, fixed points of non-expansive mappings,  $k$-smoothness,  $C^*$-algebras, $L^1$-predual space, uniform algebras.}
\begin{abstract}
 In  their seminal work, Lau and Mah (\cite{LM}) study $w^*$-normal structure in the space of operators $\mathcal{L}(H)$, on a Hilbert space $H$, using a geometric property of the dual unit ball called Lim's condition. In this paper, we study a weaker form of Lim's condition, which we call property ($\ddagger$), for $C^\ast$-algebras, uniform algebras, and $L^1$-predual spaces. In the case of a $C^\ast$-algebra, we prove that property $(\ddagger)$ is equivalent to Lim's condition and consequently, we obtain a geometric characterization of $C^*$-algebras which are $c_0$-direct sum of finite-dimensional operator spaces. For a uniform algebra, we extend a result of Lau and Mah to show that property $(\ddagger)$ implies that the space is finite-dimensional. In the case of an $L^1$-predual space, we show that this condition implies $k$-smoothness of the norm in the sense considered in \cite{LR}.
\end{abstract}

\maketitle
\section{Introduction} 
Throughout this paper, $X$ denotes a complex Banach space and $X^*$ denotes the dual of $X$. $B(X)$ and $S(X)$ are the notations used for the unit ball and the unit sphere of $X$, respectively. Geometric features of the dual unit ball of a space have strong structural consequences on the space. An interesting property in this context, which arises out of the study of $w^*$-normal structures and fixed point property, is Lim's condition. (See \cite{Lim} for the undefined notions in this context.)
\begin{Def}
  $X$ is said to satisfy Lim's condition if for every bounded net $\{\phi_{\alpha}\}$ in $X^*$ with $\|\phi_{\alpha}\|=s$, for some $s \geq 0$ and $\phi_{\alpha} \xrightarrow{w^*} 0$, we have that for any $\phi \in X^*$,
\[
\lim_{\alpha}\|\phi_{\alpha}+\phi\|=s + \|\phi\|.
\]  
\end{Def}
Note that Lau and Mah referred to this as Lim's condition on $X^*$ while we call it Lim's condition on $X$.

This property originates in the work of Lim in \cite[Theorem 2]{Lim} and is later studied by Lau and Mah in \cite{LM} as a sufficient condition for the $w^*$-normal structure and the existence of fixed points for non-expansive maps on $w^*$-compact convex sets in operator spaces. We now restate \cite[Lemma 4]{LM} as a theorem below where the authors describe weaker versions of Lim's condition.
\begin{thm}
If $X$ has Lim's condition, then
\bla
\item weak* and norm topologies coincide on the dual unit sphere $S(X^*)$ and

\item for any $0 <\e <2$, whenever $\{\phi_{\alpha}\}$ is a net in $B(X^*)$ with $\|\phi_{\alpha}-\phi_{\beta}\| \geq \e$ for all $\alpha \neq \beta$ and $\phi \in B(X^*)$ such that $\phi_{\alpha} \xrightarrow{w^\ast} \phi$, then
\[
\|\phi\| \leq 1 - \e/2.
\]
\el
\end{thm}

The authors of \cite[Theorem 2]{CPW} classify a $C^*$-algebra with weak* and weak topologies coinciding on the dual unit sphere. Hence, Lim's condition gives a classification of $C^*$-algebras via condition $(a)$ above. Thus, an interesting question is to see if condition $(b)$ above also yields a similar classification. We answer this question in the affirmative. We refer to condition $(b)$ as property $(\dag)$ for $X$.

In the geometry of Banach spaces, it is often observed that restricting a property to norm-attaining functionals is sufficient to obtain desirable results, as it contains the extremal and facial information of the space. A careful observation of the proof of \cite[Theorem 2]{CPW} shows that it is enough to assume that norm-attaining functionals are the point of weak*-weak continuity of the dual unit sphere. So, we further generalize property $(\dagger)$ by restricting the net $\{\phi_\alpha\}$ and $\{\phi\}$ to norm-attaining functionals only and call this property $(\ddagger)$.
\begin{Def}
$X$ is said to satisfy property $(\ddag)$ if for any $0 <\e <2$, whenever $\{\phi_{\alpha}\}$ is a net of norm-attaining functionals in $B(X^*)$ with $\|\phi_{\alpha}-\phi_{\beta}\| \geq \e$ for all $\alpha \neq \beta$ and $\phi \in B(X^*)$ is norm-attaining such that $\phi_{\alpha} \xrightarrow{w^\ast} \phi$, then
\[
\|\phi\| \leq 1 - \e/2.
\]
\end{Def}

Surprisingly, this apparently weaker notion of Lim's condition turns out to be equivalent to Lim's condition (hence, property $(\dagger)$) for a $C^*$- algebra. Thus, we study geometric variations of Lim's condition by considering this weaker version. One of the main results of this article proves that the property $(\ddagger)$ leads to a purely geometric characterization of $C^*$-algebras of the form
\[
\mathcal{A}=\bigoplus_{c_0}\mathcal{L}(H_\alpha),
\]
for a family $\{H_\alpha\}_{\alpha \in \Delta}$ of finite dimensional Hilbert spaces, where $\mathcal{L}(H_\alpha)$ is the space of bounded operators on $H_\alpha$. Hence, we obtain a geometric classification of the class of $C^*$-algebras which are the $c_0$-direct sum of finite-dimensional operator spaces. Thus, property $(\ddagger)$ appears to be a minimal hypothesis needed for this classification. Also, this is an improvement in the structural information on $C^*$-algebras with Lim's condition, as discussed earlier.

We also extend Lau and Mah's result that $K(\ell_2)$ does not satisfy Lim's condition to any $\ell_p$ for $1 <p <\infty$. 

As mentioned earlier, Lim's condition implies $w^*$-normal structure in a Banach space (\cite[Theorem 2]{Lim}): For every $w^*$-compact convex set $K \subset X^\ast$ and for any non-trivial convex set $H \subset K$, there is a $x_0 \in H$ such that $\sup\{\|x_0-y\|:y\in H\}< \sup\{\|x-y\|:x~,y \in H\}$. Thus, by \cite[Theorem 1]{Lim}, every non-expansive mapping  $T: K \rightarrow K$ for a $w^*$-compact convex set $K$ has a fixed point. Hence, from our classification result, the dual of $c_0$-direct sum of finite-dimensional operator spaces satisfies $w^*$-normal structure and therefore the fixed-point property holds for non-expansive maps on $w^*$-compact subsets. 

Let $\Omega$ be a compact space and let $C(\Omega)$ denote the space of continuous functions, equipped with the supremum norm. For a uniform algebra $A \subset C(\Omega)$ (i.e., a closed subalgebra containing constants and separating points), Rao in \cite[Theorem 11]{Rao} shows that if the set of points of weak*-weak continuity is weakly dense in the dual unit sphere, then A is finite-dimensional. This motivates us to analyze property $(\ddag)$ on a uniform algebra. Interestingly, we observe that any uniform algebra satisfying property $(\ddag)$ is necessarily finite-dimensional. We use similar techniques to relate property $(\ddag)$ to $k$-smoothness in the class of Banach spaces $X$ whose dual is isometric to $L^1(\mu)$ for a positive measure $\mu$. These are called $L^1$-predual or Lindenstrauss, spaces and are considered as commutative analogues of $C^\ast$-algebras. Refer to \cite[Chapter 7]{L} for this discussion. As an application, we show that property $(\ddagger)$ implies that for any  $x \in S(X)$, $x$ is a $k$-smooth point for some $k>0$ (refer to the discussion just before Theorem ~\ref{L1predual} for the definition). 

So, the main contributions of the paper are:
\bla
\item Studying a weaker version of Lim's condition in the form of property $(\ddag)$ which is shown to be equivalent to Lim's condition in $C^*$-algebras.

\item A geometric characterization of $C^*$-algebras that are $c_0$-direct sum finite-dimensional operator spaces.

\item Extension of \cite[Theorem 5]{LM} from Hilbert spaces to $\ell_p$-spaces.

\item Property $(\ddag)$ forces uniform algebras to be finite-dimensional.

\item Property $(\ddag)$ implying k-smoothness in $L^1$-predual spaces.

\el

Thus, this paper demonstrates that an apparently weaker version of Lim's condition seems to govern the structure of $C^*$-algebras, $L^1$-predual spaces, and uniform algebras in a unifying way.

This paper is organized as follows. Section 2 discusses property $(\ddag)$ in $C^*$-algebras and Section 3 is concerned with property $(\ddag)$ in function spaces.

We refer to the monographs by R. B. Holmes \cite{Ho} and J. Diestel \cite{Di} for standard results from Banach space theory, \cite{L} for isometric theory of classical Banach spaces and \cite{A1} for $C^\ast$-algebra theory.

\section{Property $(\ddag)$ for $C^\ast$ algebras}
We begin with a few basic results on property $(\ddag)$.

\begin{prop}
	If $X$ has property $(\dagger)$, then every closed subspace of $X$ has property $(\dagger)$.
\end{prop}

\begin{proof}
	Let $Y$ be a closed subspace of $X$ and $0 <\e <2$. Consider $\{\phi_\alpha\}$ to be a net in $B(Y^*)$ and $\phi \in B(Y^*)$ such that $\phi_\alpha \xrightarrow{w^*} \phi$ and $\|\phi_\alpha - \phi_\beta\| \geq \e$ for all $\alpha \neq \beta$. Let $\ol{\phi_\alpha}$ be any norm-preserving extension of $\phi_\alpha$ and assume without loss of generality that $\ol{\phi_\alpha} \xrightarrow{w^*} z^*$ for some $z^* \in B(X^*)$. Clearly, $z^*|_Y=\phi$.
	
	Notice that, we have $\| \ol{\phi_\alpha}-\ol{\phi_\beta}\| \geq \|\phi_\alpha - \phi_\beta\| \geq \e$.
	
	So, $\|\phi\| \leq \|z^*\| \leq 1-\e/2$.
\end{proof}

Similar arguments yield the following corollary.

\begin{cor} \label{corhereddag}
	If $X$ has property $(\ddag)$, then any closed subspace $Y$ of $X$ has property $(\ddag)$.
\end{cor}

We now look at property $(\ddag)$ in $\ell_p$-spaces via the following example.

\begin{ex}
Let $\{e_n\}$ be a sequence in $\ell_q$, where $e_n$ is the sequence of elements with $1$ at $n^{th}$ position and $0$ elsewhere. Let $\phi_n=e_n/(2)^{1/q}$ for all $n$ and $\phi=\phi_1$.
	
We have that, $\phi_n \xrightarrow{w^\ast} 0$. For all $n \geq 2$, let $\psi_n:=\phi_n + \phi$. Then, $\psi_n\xrightarrow{w^\ast} \phi$. Now, $\|\psi_n\|=1$ for all $n \geq 2$, that is, $\psi_n \in B(\ell_q)$. Since $\ell_q$ is reflexive, $\psi_n$ and $\phi$ are norm-attaining for each $n \in \N$. Also, for $n \neq m$, $\|\psi_n-\psi_m\|=1$.
	
But, $\|\phi\|=1/2^{1/q} >1 - 1/2$. Thus, $\ell_p$ does not satisfy property $(\ddagger)$. Also, it is known that $\ell_p$ is embedded in $\K(\ell_p)$ and by Corollary ~\ref{corhereddag}, property $(\ddag)$ is hereditary. So, $\K(\ell_p)$ does not have property $(\ddag)$.

Hence, this leads to an improvement in \cite[Theorem 5]{LM} for the space of compact operators $\K(\ell_2)$ to $\K(\ell_p)$ for $1 <p <\infty$.
\end{ex}

The result below is a substantial variation of \cite[Theorem 2]{CPW} and in particular, shows among other things that being a modular annihilator algebra can be determined by $(\ddag)$.
\begin{thm} \label{thmdagsubdif}
	
	 Let ${\mathcal A}$ be a $C^\ast$-algebra satisfying condition $(\ddag)$. Then ${\mathcal A}$ is isometrically $*$-isomorphic to  a $c_0$-direct sum $\bigoplus_{c_0} {\mathcal L}(H_{\alpha})$, for a family of finite-dimensional Hilbert spaces $\{H_{\alpha}\}_{\alpha \in J}$, for some index set $J$, and spaces of bounded operators ${\mathcal L}(H_{\alpha})$.
\end{thm}

\begin{proof}
	 We verify Condition $(iii)$ of \cite[Theorem 2]{CPW}, to get the required conclusion. Let $x \in S({\mathcal A})$  be a normal element and consider the commutative $C^*$-algebra generated by $\{x,~x^*\}$. By the Gelfand-Naimark theorem, we have the identification of this separable subalgebra as $C_0(\sigma(x))$. We also have that $\sigma(x)$ is a locally compact metric space. Since $(\ddag)$ is a hereditary property, the separable commutative $C^*$-subalgebra $C_0(\sigma(x))$ has property $(\ddag)$. By the remarks made at the start of Section 2, we get that $\sigma(x)$ is a discrete set with $0$ as the only possible accumulation point. Hence, we have the decomposition $(v)$ from \cite[Theorem 2]{CPW}. Now, by Corollary ~\ref{corhereddag} and the example following it, we see that the Hilbert spaces under consideration are finite-dimensional.
	
\end{proof}

\begin{lem} \label{lnormal}
	Let $X$ has property $(\dagger)$ and $0 \leq s \leq 1$. Then, for every $0 <\e <2s$, whenever $\{\phi_{\alpha}\}$ is a net in $X^*$ with $\|\phi_{\alpha}\| \leq s$ for all $\alpha$ and $\phi \in B(X^*)$ is such that $\phi_{\alpha} \xrightarrow{w^\ast} \phi$ and $\|\phi_{\alpha}-\phi_{\beta}\| \geq \e$ for all $\alpha \neq \beta$, then
	\[
	\|\phi\| \leq s - \e/2.
	\]
\end{lem}

\begin{proof}
	Let $0 \leq s \leq 1$ and $0 <\e <2s$. Suppose $\{\phi_{\alpha}\}$ is a net in $X^*$ with $\|\phi_{\alpha}\| \leq s$ for all $\alpha$ and $\phi \in B(X^*)$ is such that $\phi_{\alpha} \xrightarrow{w^\ast} \phi$ and $\|\phi_{\alpha}-\phi_{\beta}\| \geq \e$ for all $\alpha \neq \beta$.
	We have, $0 <\e/s <2$. Let $\psi_{\alpha}=\phi_{\alpha}/s$ for all $\alpha$ and $\psi=\phi/s$. Then, $\psi_\alpha \in B(X^*)$ for all $\alpha$, $\psi_{\alpha} \xrightarrow{w^\ast} \psi$ and $\|\psi_{\alpha}-\psi_{\beta}\| \geq \e/s$. Therefore,
	$\|\psi\| \leq 1 - \e/2s$. This implies, $\|\phi\| \leq s - \e/2$.
\end{proof}

We next consider direct sums with the maximum or supremum norm.
\begin{prop} \label{tstab}
	Let $Z = X \bigoplus_{\ell_\infty} Y$ for some Banach spaces $X$ and $Y$. Then $Z$ has property $(\dagger)$ if and only if $X$ and $Y$ have property $(\dagger)$.
\end{prop}

\begin{proof}
	Clearly, if $Z$ has property $(\dagger)$, then $X$ and $Y$ being subspaces of $Z$ and $(\dagger)$ being hereditary,	$X$ and $Y$ have property $(\dagger)$.
	
	For the proof of the converse part, we first note that $Z^* = X^* \bigoplus_{\ell_1} Y^*$.
	We denote the norms on $X$ and $Y$ also by $\|\cdot\|$.
	Let $0 <\e <2$, $\{\phi_{\alpha}\}$ be a net in $B(Z^*)$ and $\phi \in B(Z^*)$ such that $\phi_{\alpha} \xrightarrow{w^\ast} \phi$ and $\|\phi_{\alpha}-\phi_{\beta}\| \geq \e$ for all $\alpha \neq \beta$. We now choose components from the corresponding summands in the dual and also assume, without loss of generality, that the components of $\phi$ are non-zero. For each $\alpha$, let $x^*,y^*,x_{\alpha}^*, y_{\alpha}^*$ be such that $\phi_{\alpha}=x_{\alpha}^* + y_{\alpha}^*$ for all $\alpha$ and $\phi=x^* + y^*$. We have, $\|\phi_\alpha\| = \|x_{\alpha}^*\|+\|y_{\alpha}^*\| \leq 1$. Passing to subnets, if necessary, we may assume that $\|x_{\alpha}^*\|=s$ for some $0 \leq s \leq 1$ and for all $\alpha$. So, $\|y_{\alpha}^*\| \leq 1-s$. We also have, $x_\alpha^* \xrightarrow{w^\ast} x^*$ and $y_\alpha^* \xrightarrow{w^*} y^*$.
	
	We have, $\|\phi_\alpha - \phi_\beta\|=\|x_\alpha^*-x_\beta^*\| + \|y_\alpha^*-y^*\| \geq \e$. Again, passing to subnets if necessary, let $s'=\|x_\alpha^*-x_\beta^*\|$ for all $\alpha \neq \beta$. So, $\|y_\alpha^*-y_\beta^*\| \geq \e-s'$ for all $\alpha \neq \beta$.
	
	Since, $X$ and $Y$ have property $(\dagger)$, by Lemma ~\ref{lnormal},
	$\|x^*\| \leq s - (\e-s')/2$ and $\|y^*\| \leq (1-s) - s'/2$. Thus,
	\[
	\|\phi\|=\|x^*\| + \|y^*\| \leq 1 - \e.
	\]
\end{proof}

Clearly the above arguments extend to finite sums. In the proof of Proposition 9, we note that if $\phi$ attains its norm, so do the component functions. Thus Proposition 9 is also valid for components having property $(\ddag)$. 

Now, we look at the stability of property $(\ddag)$ under $c_0$-sum.
We have the following result for property $(\ddag)$ for general $c_0$-sums. Let $I$ denote the indexing set for a family of Banach spaces. For notational convenience, we omit writing the indexing set.
\begin{thm}
	Let $\{X_\gamma\}$ be any collection of  Banach spaces. Then $X=\bigoplus_{c_0}X_\gamma$ has property $(\ddag)$ if and only if $X_\gamma$ has $(\ddag)$ for each $\gamma$.
\end{thm}
\begin{proof}
Again, the hereditary nature of $(\ddag)$ implies that $X_\gamma$ has $(\ddag)$ for each $\gamma$.
	
Conversely, we have $X^*=\bigoplus_{\ell_1}X_\gamma^*$.
Again, we denote the norm on each $X_\gamma$ as $\|\cdot\|$. In what follows, we use the standard identification of $X^\ast$, in particular, when the indexing set $I$ is uncountable, $0 \neq \phi \in B(X^\ast)$ has only countably many non-zero coordinates.
	
Let $0 <\e <2$, $\{\phi_\alpha\}$ be a net of norm-attaining functionals in $B(X^*)$ with $\|\phi_\alpha - \phi_\beta\| \geq \e$ for all $\alpha \neq \beta$ and $\phi \in B(X^*)$ be a norm-attaining functional such that $\phi_{\alpha} \xrightarrow{w^\ast} \phi$.
\vskip .5em
We first note that any norm-attaining $0 \neq \phi \in B(X^*)$ has only finitely many non-zero coordinates. To see this, suppose $x \in S(X)$ and $$\|\phi\|= \sum \phi(\gamma)(x(\gamma))\leq \sum \|\phi(\gamma)\|x(\gamma)\|\leq \|\phi\|.$$
Thus, whenever $\phi(\gamma) \neq 0$, $\|x(\gamma)\|=1$ and it is easy to see that the corresponding $\phi(\gamma)$ attains its norm at $x(\gamma) \in S(X_{\gamma})$. Since $\{\|x(\gamma)\|\}$ vanishes at infinity, we obtain that $\phi$ has only finitely many non-zero coordinates.
	
	 So, without loss of generality, we may assume that for each $\alpha$, $\phi_{\alpha}$ has only first $n$ coordinates non-zero.
	
	Therefore, proceeding as in the proof of  Theorem ~\ref{tstab}, we obtain that $\|\phi\| \leq 1-\e/2$.
\end{proof}

The following is an interesting consequence of the hereditary nature of property $(\ddag)$.

\begin{cor}Let ${\mathcal A}$ be a $C^\ast$-algebra. If every separable commutative $C^\ast$-subalgebra has property $(\ddag)$, then ${\mathcal A}$ has $(\ddag)$.
\end{cor}

\begin{proof} 
It follows from the proof of Theorem ~\ref{thmdagsubdif} that the hypothesis here is sufficient to ensure the $c_0$-decomposition of ${\mathcal A}$ into finite-dimensional spaces. Thus ${\mathcal A}$ has $(\ddag)$.
\end{proof}

The following result extends part of \cite[Theorem 4]{LM} to the case of $c_0$-sums of spaces with Lim's condition. 

\begin{thm} \label{tlimsum}
	Let $\{X_i\}_{i \in \I}$ be a collection of spaces such that $X_i$ satisfies Lim's condition for each $i \in \I$. Then, $X=\bigoplus_{c_0} X_i$ satisfies Lim's condition.
\end{thm}

\begin{proof}
	The argument follows the lines of \cite[Theorem 4]{LM}. We have, $X^*=\bigoplus_{\ell_1} X_i^*$. 
      
	Let each $X_i$ satisfy Lim's condition. Let $0 \leq s \leq 1$,  $\{\phi_\alpha\}$ be a net in $X^*$ such that $\phi_\alpha \xrightarrow{w^\ast} 0$ and $\|\phi_\alpha\|=s$ for all $\alpha$. Let $\phi \in X^*$.
    
    Clearly, $\|\phi_\alpha - \phi\| \leq s + \|\phi\|$.
	
	Let $\e >0$ be given. For any finite subset $\F$ of $\I$, $\|\phi_\alpha-\phi\| \geq \|\phi_\alpha\|-\|\phi\| + 2 \sum_{i \in \F}(|\phi(i)|-|\phi_\alpha(i)|)$. 
    
    Let $\F_0$ be a finite subset of $\I$ such that $\sum_{i \in \F_0} |\phi(i)| >\|\phi\| - \e$. Since, $\phi_{\alpha} \xrightarrow{w^\ast} 0$, so there exists $\alpha_0$ such that for $\alpha \geq \alpha_0$, $\sum_{i \in \F_0} |\phi_\alpha (i)| <\e$. Thus, for $\alpha \geq \alpha_0$,
	\[
	\|\phi_\alpha - \phi\| \geq s - \|\phi\| + 2 \|\phi\| - 2\e - 2\e = s + \|\phi\| - 4\e.
	\]
	
	So, $\lim_\alpha \|\phi_\alpha - \phi\| = s + \|\phi\|$.
\end{proof}

This leads us to an interesting equivalence between property $(\ddag)$ and Lim's condition in $C^*$-algebras.
\begin{cor}
 If $\cA$ is a $C^*$-algebra with property $(\ddagger)$, then $\cA$ satisfies Lim's condition.   
\end{cor}

\begin{proof}
    Let $\cA$ satisfy property $(\ddagger)$. Then, by Theorem ~\ref{thmdagsubdif}, $\cA=\bigoplus_{c_0}\mathcal{L}(H_\alpha)$, where $H_\alpha$ is a finite-dimensional Hilbert space for each $\alpha$. 
    
    But, every finite-dimensional space satisfies Lim's condition. So, $\mathcal{L}(H_\alpha)$ satisfies Lim's condition for each $\alpha$. By the above theorem, $\cA$ satisfies Lim's condition.
\end{proof}

We recall that $Y \subset X$ is said to be a proximinal subspace if for all $x \in X$, there is a $y \in Y$ such that  $d(x,Y)= \|x-y\|$. In particular, if $P$ is a contractive projection on $X$, then $ker(P)$ is a proximinal subspace as $d(x,ker(P))=\|P(x)\|=\|x-(x-P(x))\|$. Also, if $Y\subset X$ is a subspace of finite codimension and proximinal, by Gargavi's theorem, $Y$ is a finite intersection of kernels of norm-attaining functionals. (See \cite{Sin} and \cite[Lemma 2.1]{NR}).

We do not know if Lim's condition is hereditary in general. But under some assumptions, we obtain the hereditary nature of Lim's condition.

\begin{cor}
Let $X=\bigoplus_{c_0} X_i$, where $\{X_i\}$ is a countable collection of finite-dimensional spaces. Let $Y$ be a finite co-dimensional proximinal subspace of $X$. Then, $Y$ satisfies Lim's condition.
\end{cor}

\begin{proof}
From the proof of \cite[Proposition 3.10 ((ii) $\implies$ (i))]{NR}, $Y=Y_1 \bigoplus_{\infty}Y_2$, where $Y_1$ is finite-dimensional and $Y_2=\sum_{i \in \N \setminus I}\bigoplus_{c_0}X_i$, where $I$ is a finite set. By Theorem ~\ref{tlimsum}, $Y_2$ satisfies Lim's condition. $Y_1$ being finite-dimensional also satisfies Lim's condition. Hence, $Y$ satisfies Lim's condition.
\end{proof}

The following corollary shows that Lim's condition is hereditary for $C^*$-algebras.

\begin{cor} Let $\mathcal{A}$ be a $C^\ast$-algebra. If $\mathcal{A}$ has property $(\ddag)$, then it satisfies Lim's condition and hence, it has $w^*$-normal structure as well as fixed point property for $w^*$-compact convex sets in ${\mathcal A}^\ast$.
\end{cor}

\begin{proof} 
By theorem~\ref{thmdagsubdif}, we have that $\mathcal{A}=\bigoplus_{c_0}\mathcal{L}(H_\alpha)$, where $\mathcal{L}(H_\alpha)$ is the space of bounded operators on a finite-dimensional Hilbert space $H_\alpha$. Hence, the conclusion follows from Theorem ~\ref{tlimsum}. We also get from \cite[Lemma 4]{LM} that ${\mathcal A}$ has the $w^*$-normal structure. Therefore, by \cite[Theorem 1]{Lim}, any non-expansive map $T: K \rightarrow K$ has a fixed point for every $w^*$-compact convex set $K$ in ${\mathcal A}^\ast$.
\end{proof}	

\begin{rem}
	Let $X$ satisfy Lim's condition and $Y$ be a closed subspace of $X$. Then $X/Y$ has Lim's condition. To see this, we recall that $(X/Y)^*$ is canonically identified with $Y^\perp$. Let $0 \leq s \leq 1$ and $\{\phi_\alpha\}$ be a net in $Y^\perp$ with $\|\phi_\alpha\|=s$ and $\phi_\alpha \xrightarrow{w^\ast} 0$. Let $\phi \in Y^\perp$. Since, $X$ has Lim's condition, $\lim_\alpha \|\phi_\alpha + \phi\|=s+\|\phi\|$.
\end{rem}

For preserving  property $(\ddag)$ on the quotient spaces, we need some additional hypothesis of proximinality.

\begin{prop}
	Suppose $X$ has $(\ddag)$ and $Y\subset X$ is a proximinal subspace. Then $X/Y$ has $(\ddag)$.	
\end{prop}

\begin{proof}
	We have that $Y^\perp \cong (X/Y)^*$ via the canonical isometry. Let $0 <\e <2$, $\{\phi_\alpha\}$ be a net of norm-attaining functionals in $B((X/Y)^*)$ such that $\|\phi_\alpha - \phi_\beta\| \geq \e$ and $\phi \in B((X/Y)^*)$ be a norm-attaining functional with $\phi_\alpha \xrightarrow{w^\ast} \phi$. We can identify $\{\phi_\alpha\}$ and $\phi$ as elements in $Y^\perp$.
	
	Now, $\phi$ attains norm on $B(X/Y)$. So, there exists $x_0 \in X$ such that $\phi(x_0)=\|\phi\|=d(x_0, Y)$. But, $Y$ is proximinal, thus, there is $y_0$ in $Y$ such that $\phi(x_0)=\phi(x_0-y_0)=\|\phi\|=\|x_0-y_0\|$. So, $\phi$ is norm-attaining as an element of $Y^\perp$. Similarly, for each $\alpha$, $\phi_\alpha$ is norm-attaining as an element of $Y^\perp$.
	
	Now, $(\ddag)$ is hereditary. So, $Y^\perp$ has $(\ddag)$. Thus, $(X/Y)$ has $(\ddag)$.
\end{proof}

An interesting observation relates reflexivity and Lim's condition in a dual space. If $X^\ast$ has Lim's condition, then weak$^\ast$ and weak topologies coincide on $S(X^{\ast\ast})$. By Goldstein's theorem (see \cite{Ho}), under the canonical embedding, $B(X)$ is $w^*$-dense in $B(X^{\ast\ast})$. So, we see that Lim's condition in a dual space implies reflexivity. We do not know in general if $X^\ast$ has  $(\ddag)$, then $X$ is reflexive. We have a partial result when $X$ has no isomorphic copy of $\ell_1$ ($\ell_1 \not\hookrightarrow X$).
  \begin{thm}
  	If $X$ is such that $\ell_1 \not\hookrightarrow X$ and $X^*$ has property $(\ddag)$, then $X$ is reflexive.
  \end{thm}
  \begin{proof}
  	Let $Y$ be a separable subspace of $X$ such that $\ell_1 \not\hookrightarrow Y$.  By a result in \cite[Pg. 215]{Di}, every element of $S(Y^{**})$ is a $w^*$-sequential limit of elements of $S(Y)$. Let $\phi \in S(Y^{**})$ be a norm-attaining functional and $\{y_n\}$ be a sequence in $S(Y)$ such that $y_n \xrightarrow{w^\ast} \phi$. If $\{y_n\}$ does not converge to $\phi$ in the norm, then  $\{y_n\}$ is not Cauchy. Then there is a subsequence, again denoted by $\{y_n\}$ such that for some $\e >0$, this subsequence is $\e$-separated, i.e.,  $\|y_n-y_m\| \geq \e$ for all $n \neq m$. Clearly $y_n$'s are norm-attaining on $Y^\ast$ and $y_n \rightarrow \phi$ in the $w^*$-topology.
  	
  	Now, $Y^\perp \subset X^*$ is a $w^*$-closed and hence, a proximinal subspace. Thus $Y^*$ (being isometrically isomorphic to $X^*/Y^\perp$) has property $(\ddag)$.
  	So, $\|\phi\| \leq 1-\e/2$, a contradiction. Thus, we get that $\phi \in Y$. Now by the Bishop-Phelps theorem, we conclude that $Y$ is reflexive and hence, as reflexivity is a separably determined property, $X$ is a reflexive space.
  	
  	
  \end{proof}

\section{Property $(\ddag)$ on function spaces}
For $\Omega$ locally compact, we denote by $C_0(\Omega)$ the space of functions vanishing at infinity. From the proof of \cite[Theorem 4]{LM}, we get that if $\Omega$ is a discrete set then $c_0(\Omega)$ satisfies Lim's condition. For $ \omega \in \Omega$, let $\delta(\omega)$ denote the Dirac measure at $\omega$, this is a norm-attaining functional. Note that for $ \omega_1\neq \omega_2$, $\|\delta(\omega_1)-\delta(\omega_2)\|=2$. Thus, using \cite[Theorem 4]{LM} it is easy to see that if $C_0(\Omega)$ satisfies Lim's condition or $(\ddag)$, then $\Omega$ is a discrete set.

In this section, we prove that for a uniform algebra, property $(\ddag)$ implies that the space is finite-dimensional. 

The following lemmas are crucial for the upcoming discussion.

\begin{lem} \label{lemsubalg}
Let $X \subset C(\Omega)$ be a closed subspace and suppose ${\bf 1} \in X$. Then $\overline{\partial_e B(X^\ast)}^{w^*} \subset S(X^\ast)$. All the functionals in $\overline{\partial_e B(X^\ast)}^{w^*} $ are norm-attaining.
\end{lem}

\begin{proof}
By an application of Krein-Milman theorem, any element $x^*$ of $\partial_e B(X^\ast)$ can be extended to an element of $\partial_e C(\Omega)^\ast$. So, $|x^*({\bf 1})|=1$. Now, $w^*$-limit of such a point also takes absolute value $1$ at the constant function ${\bf 1}$. Thus,  $\overline{\partial_e B(X^\ast)}^{w^*} \subset S(X^\ast)$. All the functionals in the $w^*$-closure attain their norm at ${\bf 1}$.
\end{proof}

The following lemma appears as \cite[Thoerem 2]{IL}. We prove it for the sake of completeness.

\begin{lem} \label{lemonedsum}
	Let $x^* \in \partial_e B(X^*)$ be a one dimensional $L$-summand, that is, $X^*=span\{x^*\} \bigoplus_1 N$ for some closed subspace $N$ of $X^*$. Then, $x^*$ attains its norm.
\end{lem}

\begin{proof}
	Since $X^\ast = span\{x^\ast\} \bigoplus_1 N$, we note that if $y^\ast \in S(X^\ast)$ is norm-attaining and $y^\ast = \alpha x^\ast+n^\ast$ for some $n^\ast \in N$, with both the components non-zero, then it is easy to see that $x^\ast$ also attains its norm. Now, using the Bishop-Phelps theorem on the denseness of norm-attaining functionals (\cite[Pg. 169]{Ho}), we get that there is a norm-attaining functional $y^\ast$ for which the components in the decomposition are non-zero. Thus $x^\ast$ attains its norm.
\end{proof}

  \vskip 1em
   We recall that if $X$ is an $L^1$-predual space, then since $X^\ast = L^1(\mu)$, for any $x^\ast \in \partial_e B(X^*)$ we have $X^\ast = span\{x^\ast\} \bigoplus_1 N$. So, by Lemma ~\ref{lemsubalg}, $x^*$ attains norm. This is the motivation for studying $(\ddag)$ in this context.
  \vskip 1em

In the case of a uniform algebra $A$, it follows from the properties of the Choquet boundary of $A$, that for any extreme point $x^*$ of the dual unit ball, $span\{x^*\}$ is an $L$-summand (see Theorem V.4.2 in \cite{HWW} and Lemma II.12.2 in \cite{Ga}). For a compact convex set $K$, let $A(K)$ denote the space of affine continuous functions. For an extreme point $k$ of $K$, that is also a split face, it is known that $span(\delta(k))$ is a one dimensional $L$-summand in $A(K)^*$ (see \cite[Chapter 3, Theorem 8.3]{AE}). This in particular happens when $K$ is a Choquet Simplex. To connect it with the $C^*$ algebra theory, it follows from \cite[Thoerem 1.1]{BR} that the set of tracial state $T(\mathcal{A})$ for a unital $C^*$-algebra $\mathcal{A}$ is a Choquet simplex. These spaces illustrate the significance of the next set of results.

\begin{thm} \label{thmunialg}
Let $X$ be a Banach space such that the set $D= \{x^\ast \in \partial_e B(X^*): span\{x^\ast\}~is~an~L-summand\}$ is $w^*$-dense in $\partial_e B(X^*)$. Suppose $\overline{\partial_e B(X^*)}^{w\ast} \subset S(X^*)$. If $X$ has $(\dag)$, then $X$ is isometric to $\ell^\infty(k)$ for some integer $k>0$. In particular, any uniform algebra satisfying $(\ddag)$ is finite-dimensional.
\end{thm}

\begin{proof} 

Let $x^\ast$ be a $w^*$-accumulation point of $D$. Let $\{x^\ast_{\alpha}\} \subset D$ be a net such that $x^\ast_{\alpha} \rightarrow x^\ast$ in the $w^*$ topology. Since any two independent  points in $D$ can not belong to the same one dimensional $L$-summand, we may and do assume that $\|x^\ast_{\alpha}-x^\ast_{\beta}\| =2 $ for $\alpha \neq \beta$. Since $\overline {\partial_e B(X^*)}^{w^*} \subset S(X^\ast)$, we get $\|x^\ast\|=1$. This contradicts property $(\dag)$.

This shows that $D$, after identifying vectors that are multiples of scalars of modulus $1$, is a finite  discrete set. Thus $X$ is isometric to $\ell^\infty(k)$. In the case of a uniform algebra, we note, by Lemma ~\ref{lemsubalg}, that extreme points of the dual unit ball as well as their $w^*$-accumulation points are norm-attaining functionals.	Therefore, any uniform algebra satisfying property $(\ddag)$ is finite-dimensional.
\end{proof}

As a consequence of the above result, we have

\begin{cor}
If $C(\Omega)$ has property $(\ddag)$, then $\Omega$ is finite.
\end{cor}

Let $X$ be a Banach space such that $X^\ast$ is isometric to $L^1(\mu)$ for some positive measure $\mu$. We recall that for any compact or locally compact set $\Omega$, the spaces $C(\Omega)$ and $C_0(\Omega)$, $A(K)$-spaces for a simplex $K$ are in this class (see \cite{L}). 

Recall from \cite{LR} that a non-zero element $x \in X$ is said to be a \textit{k-smooth} point if the set of supporting functionals $J(x):=\{f \in S(X^*): f(x)=\|x\|\}$ is $k$-dimensional, that is, $dim(span(J(x))=k$.

\begin{thm} \label{L1predual}
Let $X$ be an $L^1$-predual space with $(\ddag)$.  For any $x \in S(X)$, $x$ is a $k$-smooth point for some $k>0$.
 \end{thm}
 
 \begin{proof} Let $x_0 \in S(X)$. Denote by $S_{x_0}:=\{f \in S(X^*) : f(x_0)=1\}$. Since $X$ is an $L^1$-predual space and as $S_{x_0}$ is a w$^*$-closed extreme convex subset (face) of $B(X^*)$,
 it follows from \cite[Theorem 3.3]{E}  that $span\{S_{x_0}\}$ is a w$^*$-closed subspace and $X^\ast = span\{S_{x_0}\} \bigoplus_1 N$ for a closed subspace $N \subset X^\ast$. 
 
 Now, if $J = \{x \in X: x(S_{x_0})=0\}$, then $J^\bot = span\{S_{x_0}\}$. Let $\pi: X \rightarrow X/J$ be the quotient map and  $\Phi: X/J \rightarrow A(S_{x_0})$ be defined as $\Phi(x)(x^\ast)= x^\ast(x_0)$, where $x^\ast \in S_{x_0}$ and $x \in X$. This is a linear map taking $\pi(x_0)$ to $1$ and it is easy to see that $\Phi$ is a surjective isometry. 
 
 It now follows from \cite[Theorem 2, Section 19]{L} and its complex versions from \cite[Section 20]{L}, that $S_{x_0}$ is a Choquet simplex. Since $X/J$ has $(\ddag)$, we obtain from Theorem~\ref{thmunialg} that $X/J$ is isometric to $\ell^\infty(k)$. Consequently $x$ is a $k$-smooth point.
 \end{proof}
 
  \begin{rem} It is reasonable to conjecture that if $X$ is a $L^1$-predual space and has $(\ddag)$, then it is isometric to $c_0(\Gamma)$ for a discrete set $\Gamma$.
\end{rem}

 \begin{rem}
In \cite{CPW}, the authors characterize Kaplansky's $I$-rings among $C^\ast$-algebras in terms of denseness of points of strong subdifferentiability in the unit sphere. It would be interesting	to get an analogue of this in terms of variations on Lim's conditions or fixed point property that we have considered here.  Classification schemes related to these ideas will appear elsewhere.
 \end{rem}
\smallskip
Some of these results are part of an ongoing research of the second author, funded by the Anusandhan National Research Foundation (ANRF) Core Research Grant 2024-2027, CRG-2023/ No-000595, `Classification of Banach spaces using differentiability'.
\smallskip
\begin{thebibliography}{99}
\bibitem{A1} W. B. Arveson, {\em An invitation to $C^*$-algebras}, Springer GTM 39, Berlin 1976.

\bibitem{AE} L. Asimow and A.~J. Ellis, {\em Convexity theory and its applications in functional analysis}, London Mathematical Society Monographs, 16, Academic Press, London-New York, 1980.

\bibitem{B}S. Basu and T. S. S. R. K. Rao, {\em Some stability results for asymptotic norming properties in Banach spaces}, Colloq. Math. 75 (1998) 271--284.

\bibitem{BR} B. Blackader and M. R{\o}rdam,{\em The space of tracial states on a $C^*$-algebra},  Expositiones Mathematicae (2024) 125618,
https://doi.org/10.1016/j.exmath.2024.125618. 

\bibitem{CPW} M. D. Contreras, R. Pay\'a and W. Werner, {\em $C^*$-algebras that are I-rings}, J. Math. Anal. Appl. 198 (1996) 227--236.

\bibitem{Di} J. Diestel, \emph{Sequences and series in Banach spaces}, Graduate Texts in Mathematics, 92, Springer, New York, 1984.

\bibitem{E} A. J. Ellis, A. K. Roy, T. S. S. R. K. Rao and U. Uttersrud, {\em Facial structure of complex Lindenstrauss spaces}, Trans. Amer. math. Soc., 268 (1981) 173--186.
    
\bibitem{Ga} T. Gamelin, {\em Uniform algebras}, Prentice-Hall, Englewood Cliffs, 1969.
    
\bibitem{HWW} P. Harmand, D. Werner and W. Werner, {\em $M$-ideals in Banach spaces and Banach algebras}, Springer LNM 1547, Berlin 1993.

\bibitem{Ho} R. B. Holmes, {\em Geometric Functional analysis and its applications}, Springer. GTM 24, Berlin 1975.

\bibitem{IL} V. Indumathi and S. Lalithambigai, \emph{A new proof of proximinality for $M$-ideals}, Proc. Amer. Math. Soc. {\bf 135} (2007), no.~4, 1159--1162.

\bibitem{L} H. E.  Lacey, {\em  The isometric theory of classical Banach spaces}, In: DieGrundlehren der Mathematischen Wissenschaften, Band, vol. 208, pp. x+270. Springer, New York (1974).

\bibitem{Lim} T. C. Lim, {\em Asymptotic centers and nonexpansive mappings in conjugate spaces}, Pacific J. Math., 90 (1980) 135--143.

\bibitem{LM} A. T. M. Lau and P. F. Mah, {\em Quasi-normal structure for certain spaces of operators on Hilbert spaces}. Pacific J. Math. 121 (1986) 109--118.

\bibitem{LR} B. L. Lin and T. S. S. R. K. Rao, {\em Multi smoothness in Banach spaces}, Int. J. Math. Math. sci. (2007) 52382  Art.ID , 12pp.

\bibitem{NP} I. Namioka and R. R. Phelps, {\em Banach spaces which are Asplund spaces}, Duke Math. J., 42 (1975) 735--750.

\bibitem{NR} D. Narayana and T.~S.~S.~R.~K. Rao, {\em Transitivity of proximinality and norm-attaining functionals}, Colloq. Math., 104 (2006),
 1--19.

\bibitem{Rao} T.S.S.R.K. Rao, \emph{Uniform algebras as Banach spaces} Complex Anal. Oper. Theory, (2026), article 48. 
 
 \bibitem{Sin} I. Singer, {\em Best approximations in normed linear spaces by elements of linear subspaces}, Springer, Berlin, 1970.
\end{thebibliography}
\end{document}

